\subsection*{不要先追两根，先追它们的中心}
若两个数的和已知，就知道它们的平均值。把两根写成 $m-d,m+d$，和自动成为 $2m$，只剩一个未知偏移 $d$。它们的积是 $m^2-d^2$，因此给定和与积之后，问题只剩求一个平方。

对 $x^2-10x+21=0$，两根中心是五，偏移平方是 $25-21=4$，于是根为三和七。这里不是省略了求根原理，而是选择了让对称性直接显现的变量。配方 $(x-5)^2=4$ 是同一想法的另一种写法。

\subsection*{从特殊形式回到一般二次式}
面对 $Ax^2+Bx+C=0$，先确认 $A\ne0$，再除以它，把二次项系数变为一。中心是 $-B/(2A)$，剩下的判别信息来自偏移平方。若它为零，两根重合；若为负，实数中无根，但复数中仍可解释。

还可以从抛物线看：平移到对称轴后，一次项消失，左右两交点围绕中心成对出现。方程的代数对称性和图像对称轴于是对应起来。

解完后检查两根的和与积，能快速发现符号错误。这个方法以后还会出现在二维特征值计算中：迹给根的和，行列式给根的积。基础代数不是独立的技巧仓库，它会在更高层问题里反复提供结构。
